% oge_ref.tex
% Расширенные справочные данные ОГЭ по физике (9 класс)
\begingroup
\footnotesize
\setlength{\parskip}{2pt}
\setlength{\tabcolsep}{3.5pt}
\renewcommand{\arraystretch}{0.92}

\subsection*{Десятичные приставки}
\noindent\begin{tabularx}{\textwidth}{|>{\raggedright\arraybackslash}X|c|c||>{\raggedright\arraybackslash}X|c|c|}
\hline
\textbf{Приставка} & \textbf{Обозначение} & \textbf{Множитель} & \textbf{Приставка} & \textbf{Обозначение} & \textbf{Множитель} \\ \hline
мега  & М & $10^6$  & санти & с & $10^{-2}$ \\ \hline
кило  & к & $10^3$  & милли & м & $10^{-3}$ \\ \hline
гекто & г & $10^2$  & микро & мк & $10^{-6}$ \\ \hline
деци  & д & $10^{-1}$ & нано  & н & $10^{-9}$ \\ \hline
\end{tabularx}

\subsection*{Константы и справочные величины}
\noindent\begin{tabularx}{\linewidth}{>{\raggedright\arraybackslash}X r}
Ускорение свободного падения на Земле & $g = 10~\text{м/с}^2$ \\
Скорость звука в воздухе (при $20~^\circ\text{C}$) & $v_{\text{зв}} = 340~\text{м/с}$ \\
Скорость света в вакууме & $c = 3{,}0 \cdot 10^8~\text{м/с}$ \\
Элементарный электрический заряд & $e = 1{,}6 \cdot 10^{-19}~\text{Кл}$ \\
Нормальное атмосферное давление & $p_0 = 10^5~\text{Па} = 760~\text{мм рт. ст.}$ \\
\end{tabularx}

\subsection*{Плотности веществ}
\noindent\begin{tabularx}{\linewidth}{X r || X r}
\hline
Вода & $1000~\text{кг/м}^3$ & Керосин & $800~\text{кг/м}^3$ \\ \hline
Лёд & $900~\text{кг/м}^3$   & Ртуть & $13600~\text{кг/м}^3$ \\ \hline
Воздух & $1{,}29~\text{кг/м}^3$ & Сталь, железо & $7800~\text{кг/м}^3$ \\ \hline
Алюминий & $2700~\text{кг/м}^3$ & Медь & $8900~\text{кг/м}^3$ \\ \hline
\end{tabularx}

\subsection*{Удельная теплоёмкость}
\noindent\begin{tabularx}{\linewidth}{X r || X r}
\hline
Вода & $4200~\text{Дж/(кг}\cdot{^\circ}\text{C)}$ & Лёд & $2100~\text{Дж/(кг}\cdot{^\circ}\text{C)}$ \\ \hline
Алюминий & $920~\text{Дж/(кг}\cdot{^\circ}\text{C)}$  & Сталь & $500~\text{Дж/(кг}\cdot{^\circ}\text{C)}$ \\ \hline
Медь & $400~\text{Дж/(кг}\cdot{^\circ}\text{C)}$  & Свинец & $130~\text{Дж/(кг}\cdot{^\circ}\text{C)}$ \\ \hline
\end{tabularx}

\subsection*{Удельная теплота плавления и парообразования}
\noindent\begin{tabularx}{\linewidth}{X r || X r}
\hline
Плавление льда & $3{,}3 \cdot 10^5~\text{Дж/кг}$ & Парообразование воды & $2{,}3 \cdot 10^6~\text{Дж/кг}$ \\ \hline
Плавление свинца & $2{,}5 \cdot 10^4~\text{Дж/кг}$ & Сгорание бензина & $4{,}6 \cdot 10^7~\text{Дж/кг}$ \\ \hline
\end{tabularx}

\subsection*{Температуры плавления и кипения веществ}
\noindent\begin{tabularx}{\linewidth}{X r || X r}
\hline
Плавление льда & $0~^\circ\text{C}$ & Кипение воды & $100~^\circ\text{C}$ \\ \hline
Плавление свинца & $327~^\circ\text{C}$ & Кипение спирта & $78~^\circ\text{C}$ \\ \hline
Плавление меди & $1085~^\circ\text{C}$ & Кипение ртути & $357~^\circ\text{C}$ \\ \hline
\end{tabularx}

\subsection*{Удельное электрическое сопротивление проводников (при $20~^\circ\text{C}$)}
\noindent\begin{tabularx}{\linewidth}{X r || X r}
\hline
Серебро & $0{,}016~\frac{\text{Ом}\cdot\text{мм}^2}{\text{м}}$ & Железо & $0{,}10~\frac{\text{Ом}\cdot\text{мм}^2}{\text{м}}$ \\ \hline
Медь    & $0{,}017~\frac{\text{Ом}\cdot\text{мм}^2}{\text{м}}$ & Никелин & $0{,}40~\frac{\text{Ом}\cdot\text{мм}^2}{\text{м}}$ \\ \hline
Алюминий & $0{,}028~\frac{\text{Ом}\cdot\text{мм}^2}{\text{м}}$ & Нихром & $1{,}1~\frac{\text{Ом}\cdot\text{мм}^2}{\text{м}}$ \\ \hline
\end{tabularx}
\vspace{6pt}
\endgroup
\newpage
